\documentclass[11pt]{article}

\usepackage{amsmath, amssymb}

\title{Density Behavior Under Primorial Sieve and Structured Modular Constraints}
\author{}
\date{}

\begin{document}

\maketitle

\section{Domain Definition}

Let
\[
S_L = \{ n \leq L : n \equiv 5 \pmod{6} \}.
\]

We define the normalized density:
\[
D(L) = \frac{|S_L|}{L/6}.
\]

All results are evaluated within this progression.

\section{Primorial Sieve}

\subsection{Definition}

Let $p_1, p_2, \dots, p_k$ denote the first $k$ primes and define the primorial:
\[
P_k = \prod_{i=1}^k p_i.
\]

We consider the filtered set:
\[
S_{k,L} = \{ n \leq L : n \equiv 5 \pmod{6}, \ \gcd(n, P_k) = 1 \}.
\]

\subsection{Reduction}

Since $n \equiv 5 \pmod{6}$ implies $\gcd(n,6)=1$, primes $2$ and $3$ do not contribute additional filtering.

Thus, only primes $p \geq 5$ affect the density.

\subsection{Density Formula}

Assuming independence of residue classes modulo distinct primes, the density is:
\[
D_k = \prod_{\substack{p \mid P_k \\ p \geq 5}} \left(1 - \frac{1}{p}\right).
\]

\subsection{Asymptotic Behavior}

As $k \to \infty$,
\[
\prod_{p \geq 5} \left(1 - \frac{1}{p}\right) = 0,
\]
by divergence of the harmonic series over primes.

Therefore:
\[
D_k \to 0.
\]

\section{Structured Modular Constraint}

\subsection{Definition}

Let $m$ be a positive integer such that $\gcd(m,6)=1$.

Define:
\[
S_{m,L} = \{ n \leq L : n \equiv 5 \pmod{6}, \ n \not\equiv 0 \pmod{m} \}.
\]

\subsection{Density Formula}

Within the progression $n \equiv 5 \pmod{6}$, residue classes modulo $m$ are uniformly distributed when $\gcd(m,6)=1$.

Thus:
\[
D = 1 - \frac{1}{m}.
\]

\subsection{Example}

For $m = 25$:
\[
D = 1 - \frac{1}{25} = \frac{24}{25}.
\]

\section{Multi-Constraint System}

\subsection{Definition}

Let $m_1, \dots, m_r$ be pairwise coprime integers with $\gcd(m_i,6)=1$.

Define:
\[
S_{L} = \{ n \leq L : n \equiv 5 \pmod{6}, \ n \not\equiv 0 \pmod{m_i} \ \forall i \}.
\]

\subsection{Density Formula}

By independence of residue classes modulo pairwise coprime moduli:
\[
D = \prod_{i=1}^r \left(1 - \frac{1}{m_i}\right).
\]

\subsection{Example}

For $m_1 = 25$, $m_2 = 49$:
\[
D = \left(1 - \frac{1}{25}\right)\left(1 - \frac{1}{49}\right)
= \frac{24}{25} \cdot \frac{48}{49}.
\]

\section{Comparison of Systems}

\subsection{Primorial Sieve}

\[
D_k = \prod_{p \geq 5} \left(1 - \frac{1}{p}\right) \to 0.
\]

\subsection{Structured Constraints}

\[
D = \prod_{i=1}^r \left(1 - \frac{1}{m_i}\right),
\]
which remains bounded for finite $r$.

\subsection{Conclusion}

The two systems exhibit fundamentally different behavior:

\begin{itemize}
\item Multiplicative prime filtering produces density decay.
\item Finite modular constraints preserve bounded density.
\end{itemize}

Therefore, density collapse is not a universal property of filtering, but a consequence of unbounded multiplicative accumulation.

\section{Remarks}

\begin{itemize}
\item All formulas assume uniform distribution across admissible residue classes.
\item Independence holds exactly for coprime moduli via the Chinese Remainder Theorem.
\item The distinction between infinite prime products and finite modular products is structural, not numerical.
\end{itemize}

\end{document}
